\subsection{次数}

考虑一个环 B 与 B 的一个子环 A。给 B 赋以由 \(_{s}B_{d}\) 上的 (B, B)-双模结构经标量限制导出的 (A, A)-双模结构。

命题 5. —— 设 B 是一个环，A 是 B 的一个单子环，S 是一个单左 A-模。那么 B 是一个自由左 A-模，其维数为 \(\text{long}_{\text{A}}(\text{B} \otimes_{\text{A}} \text{S})\)。

设 r 是 A 的长度；A-模 \(A_{s}\) 同构于 \(S^{r}\)。而左 A-模 B 同构于 \(B \otimes_{A} A_{s}\)（II, §3, No.~4, 第 249 页），因而同构于 \((\mathrm{B} \otimes_{\mathrm{A}} \mathrm{S})^{r}\)（II, §3, No.~7, 第 255 页，命题 7）。于是我们有 \(\operatorname{long}_{\mathrm{A}}(\mathrm{B}) = r \operatorname{long}_{\mathrm{A}}(\mathrm{B} \otimes_{\mathrm{A}} \mathrm{S})\)，从而命题 5 由 VIII, 第 123 页的命题 3 得出。

定义 2. —— 设 B 是一个环，A 是 B 的一个单子环。自由左 A-模 B 的维数称为 B 在 A 上的（左）次数，记作 \(^{(1)}\) \([B:A]_{s}\)。

把 A 与 B 换成它们的反环，由上所述可推出 B 是一个自由右 A-模。我们用 \([B : A]_{d}\) 表示它的维数，并称之为 B 在 A 上的右次数。

\pandocbounded{\includegraphics[keepaspectratio]{images/part001_1658b52d.jpg}}

我们可以给出一个体 B 与一个子体 A 的例子，使得次数 \([B:A]_{s}\) 与 \([B:A]_{d}\) 不同。\(^{(2)}\)

设 B 是一个环，A 是 B 的一个单子环，S 是一个单左 A-模。设 M 是一个左 A-模，a 是它的长度。A-模 M 与 \(S^{(a)}\) 同构（VIII, 第 122 页，命题 2），因而 B-模 \(B \otimes_{A} M\) 与 \((\mathrm{B} \otimes_{\mathrm{A}} \mathrm{S})^{(a)}\) 亦然。关系式

\[
\mathrm{long} _ {\mathrm{A}} (\mathrm{B} \otimes_ {\mathrm{A}} \mathrm{M}) = [ \mathrm{B}: \mathrm{A} ] _ {s} \mathrm{long} _ {\mathrm{A}} (\mathrm{M})\tag{3}
\]

由命题 5 与定义 2 得出。

命题 6. —— 设 C 是一个环，B 是 C 的一个单子环，A 是 B 的一个单子环。那么我们有 \([C:A]_{s} = [C:B]_{s}[B:A]_{s}\)。

引入把 C 视为左 B-模时的一个基 \((e_{i})_{i\in\mathrm{I}}\)，以及把 B 视为左 A-模时的一个基 \((f_{j})_{j\in\mathrm{J}}\)。那么族 \((f_{j}e_{i})_{j\in\mathrm{J},i\in\mathrm{I}}\) 是把 C 视为左 A-模时的一个基（II, §1, No.~13, 第 222 页，命题 25）；命题 6 得证。

注. —— 1) 设 A 是单环 B 的一个单子环，且右次数 \([B : A]_{d}\) 有限。设 C 是把 B 视为右 A-模时的自同态环；由 VIII, 第 123 页的命题 4, b)，它是一个单环。对 B 中的任意 b，设 \(\gamma(b)\) 是从 B 到 B 的映射 \(x \mapsto bx\)；则 \(\gamma : b \mapsto \gamma(b)\) 是从 B 到 C 的一个子环的同构。此外，若 \((x_{1}, \ldots, x_{m})\) 是右 A-模 B 的一个基，则把 c 映为 \((c(x_{1}), \ldots, c(x_{m}))\) 的态射是从 C 到 \(B_{s}^{m}\) 的左 B-模同构；于是我们有关系式

\[
[ \mathrm{C}: \gamma (\mathrm{B}) ] _ {s} = [ \mathrm{B}: \mathrm{A} ] _ {d}.\tag{4}
\]

把 VIII, 第 123 页的关系式 (2) 应用于右 A-模 B，我们看出

\[
\mathrm{long} (\mathrm{C}) = [ \mathrm{B}: \mathrm{A} ] _ {d} \mathrm{long} (\mathrm{A}).\tag{5}
\]

\begin{enumerate}
\def\labelenumi{\arabic{enumi})}
\setcounter{enumi}{1}
\tightlist
\item
  设 \(\mathrm{K}\) 是一个交换域（域）。若 \(\mathrm{A}\) 是有限次数代数 \(\mathrm{B}\)（在 \(\mathrm{K}\) 上）的一个单子代数，则由 VIII, 第 125 页的命题 6，\(\mathrm{B}\) 在 \(\mathrm{A}\) 上的左次数满足关系式 \(\left[\mathrm{B}:\mathrm{A}\right]_s[\mathrm{A}:\mathrm{K}] = [\mathrm{B}:\mathrm{K}]\)。同样，我们有 \(\left[\mathrm{B}:\mathrm{A}\right]_d[\mathrm{A}:\mathrm{K}] = [\mathrm{B}:\mathrm{K}]\)。由此得等式 \(\left[\mathrm{B}:\mathrm{A}\right]_s = \left[\mathrm{B}:\mathrm{A}\right]_d\)。
\end{enumerate}

